% Journal of Psychoactive Drugs — submission-format skeleton
% Built from verified requirements in ../JOURNAL_REQUIREMENTS.md (no official T&F LaTeX class
% for this journal was available; JPD is not in the venue-templates skill's bundled set).
% Structure only — Phase 3 fills this in. Do not add content directly here; edit drafts/v1_draft.tex.
%
% Verified formatting (see ../JOURNAL_REQUIREMENTS.md for sourcing/confidence per line):
%   - 4,000 words max, Introduction–Methods–Results–Discussion combined (verbatim from JPD instructions)
%   - Unstructured abstract, <=200 words (verbatim)
%   - Chicago Manual of Style, Author-Date references (verbatim) — NOT APA
%   - Double-spaced, 1in margins, numbered pages, 12pt Times New Roman or similar, American spelling
%   - 4-6 keywords

\documentclass[12pt]{article}

\usepackage[utf8]{inputenc}
\usepackage[T1]{fontenc}
\usepackage{times}                 % 12pt Times New Roman per JPD spec
\usepackage[margin=1in]{geometry}  % 1-inch margins per JPD spec
\usepackage{setspace}
\doublespacing                     % double-spaced per JPD spec
\usepackage{lineno}                % line numbers — standard for peer review, harmless if unrequired
\usepackage{graphicx}
\usepackage{booktabs}
\usepackage{amsmath}

% --- References: Chicago Author-Date ---
% Primary choice: natbib + chicago.bst (classic "Chicago" BibTeX style, author-date variant).
% Confirm chicago.bst is available wherever compilation happens (see JOURNAL_REQUIREMENTS.md —
% not found on this machine's PATH; not a Phase 0 blocker, must be confirmed before Phase 6).
\usepackage[authoryear,round]{natbib}
\bibliographystyle{chicago}

% --- Alternative (uncomment if biber/biblatex-chicago is available instead) ---
% \usepackage[authordate,backend=biber,natbib=true]{biblatex-chicago}
% \addbibresource{../references/references.bib}

\title{[TITLE — descriptive, first-of-kind framing per MANUSCRIPT_PLAN.md Phase decisions]}

\author{
  [Author One]$^{1}$ \and [Author Two]$^{2}$ \and \ldots \\
  \small $^{1}$[Affiliation One] \\
  \small $^{2}$[Affiliation Two] \\
  \small Corresponding author: [name, email] % author order/affiliations — open item, see MANUSCRIPT_PLAN.md
}

\date{}

\begin{document}

\maketitle
\linenumbers

\begin{abstract}
\noindent
% UNSTRUCTURED, <=200 words. No Background/Methods/Results/Conclusions headers — JPD requirement,
% confirmed against 5 recently published JPD survey articles with zero exceptions.
% Written LAST, after Phases 0-5 are complete and every number is final.
[abstract text]
\end{abstract}

\noindent\textbf{Keywords:} [keyword 1]; [keyword 2]; [keyword 3]; [keyword 4] % 4-6 required

\section{Introduction}
% Phase 2 citations land here. Target: keep this section lean — the 4,000-word cap covers
% Intro+Methods+Results+Discussion combined, and the current draft is already over budget.
[content]

\section{Methods}
\subsection{Design}
[content]
\subsection{Recruitment and sample construction}
% Phase 1d: recruitment/sampling facts recovered from the proposal narrative go here —
% do NOT leave as [not reported in source], that gap is closed.
[content]
\subsection{Instrument and variable definitions}
% Phase 1c: verbatim item wording + scale anchors from koboxls.xlsx
[content]
\subsection{Missing data, structural skips, and suppression}
[content]
\subsection{Statistical approach}
[content]
\subsection{Protocol deviations}
[content]

\section{Results}
% Reuse draft's six subsections; patch every number resolved in Phase 1.
[content]

\section{Discussion}
% Substantially expanded vs. current draft using Phase 2c comparator literature —
% this is the section most likely to push the word count over budget; watch it closely.
[content]

\section{Limitations}
[content]

\section{Conclusion}
[content]

\section*{Author contributions}
% CRediT roles — required disclosure per JOURNAL_REQUIREMENTS.md
[content]

\section*{Funding}
[content]

\section*{Competing interests}
[content]

\section*{Data availability}
[content]

\section*{Generative AI use declaration}
% Required JPD disclosure — state LLM-assisted drafting/analysis tooling used in this project.
[content]

\bibliography{../references/references}

% --- Supplementary material (NOT counted against the 4,000-word cap) ---
% Everything the word-budget note in JOURNAL_REQUIREMENTS.md moves out of the main text:
% full BH-correction-by-family table, facilitator self-report detail table, ordinal-model table,
% STROBE checklist, instrument documentation table, qualitative coding scaffold.
% Build as a separate supplement.tex, not appended here.

\end{document}
